\documentclass[11pt]{article}
\usepackage[margin=0.85in]{geometry}
\usepackage{amsmath,amssymb}
\usepackage{enumitem}
\usepackage{titlesec}
\usepackage{xcolor}
\usepackage{mdframed}
\usepackage{hyperref}
\hypersetup{colorlinks=true,linkcolor=blue,linktoc=all}

\titleformat{\section}{\normalfont\Large\bfseries}{}{0em}{}
\titleformat{\subsection}{\normalfont\bfseries\color{blue!60!black}}{Q\thesubsection}{0.5em}{}
\renewcommand{\thesubsection}{\arabic{subsection}}

\newmdenv[
  linewidth=0.6pt,
  linecolor=gray!60,
  backgroundcolor=gray!5,
  innermargin=8pt,
  outermargin=0pt,
  skipabove=6pt,
  skipbelow=10pt
]{qbox}

\newcommand{\Sol}{\vspace{2pt}\noindent\textbf{Solution.}\ }

\setlength{\parindent}{0pt}
\setlength{\parskip}{2pt}
\allowdisplaybreaks

\title{\vspace{-2cm}Stat 110 --- Chapter 1: Probability and Counting\\[2pt]\large 35 Problems with Full Step-by-Step Solutions}
\author{}
\date{}

\begin{document}
\maketitle
\vspace{-1cm}
\tableofcontents
\newpage

\section*{Textbook Exercises}
\addcontentsline{toc}{section}{Textbook Exercises}

%============================================================
\subsection{Book Ch.1, Exercise 8}
\begin{qbox}
(a) How many ways are there to split a dozen people into 3 teams, where one team has 2 people, and the other two teams have 5 people each?

(b) How many ways are there to split a dozen people into 3 teams, where each team has 4 people?
\end{qbox}
\Sol
\textbf{(a)} Pick 2 of the 12 for the 2-person team, then 5 of the remaining 10 for a team of 5; the last 5 form the other team. This overcounts by $2$ since the two 5-person teams are not distinguished.
\begin{align*}
\binom{12}{2}\binom{10}{5} &= 66 \times 252 = 16632\\
\frac{16632}{2} &= 8316
\end{align*}
\textbf{(b)} Line up all 12 people ($12!$ ways); the first 4 form Team A, next 4 Team B, last 4 Team C. Order within a team doesn't matter ($4!$ each), nor does which team is "first," "second," "third" ($3!$ orderings of the teams):
\begin{align*}
12! &= 479001600\\
4!\cdot 4!\cdot 4!\cdot 3! &= 24\times 24\times 24\times 6 = 82944\\
\frac{479001600}{82944} &= 5775
\end{align*}

%============================================================
\subsection{Book Ch.1, Exercise 9}
\begin{qbox}
(a) How many paths are there from the point $(0,0)$ to the point $(110,111)$ in the plane such that each step either consists of going one unit up or one unit to the right?

(b) How many paths are there from $(0,0)$ to $(210,211)$, where each step consists of going one unit up or one unit to the right, and the path has to go through $(110,111)$?
\end{qbox}
\Sol
\textbf{(a)} Encode a path as a string of $U$'s and $R$'s. To reach $(110,111)$ we need exactly $110$ R's and $111$ U's, total length $221$. The path is determined by choosing which $110$ of the $221$ positions are R's:
\begin{align*}
\binom{110+111}{110} = \binom{221}{110}
\end{align*}
\textbf{(b)} By the multiplication rule: (paths to $(110,111)$) $\times$ (paths from $(110,111)$ to $(210,211)$). The second leg needs $210-110=100$ R's and $211-111=100$ U's:
\begin{align*}
\binom{221}{110}\times\binom{100+100}{100} = \binom{221}{110}\binom{200}{100}
\end{align*}

%============================================================
\subsection{Book Ch.1, Exercise 15}
\begin{qbox}
Give a story proof that $\displaystyle\sum_{k=0}^n \binom{n}{k}=2^n$.
\end{qbox}
\Sol
Consider picking a subset of $n$ people.
\begin{itemize}[leftmargin=1.3em]
\item \emph{Counting by size:} for each $k=0,1,\dots,n$, there are $\binom{n}{k}$ subsets of size exactly $k$. Summing over all sizes counts every subset exactly once:
\begin{align*}
\#\{\text{subsets}\} = \sum_{k=0}^n \binom{n}{k}
\end{align*}
\item \emph{Counting directly:} each of the $n$ people is independently either in or out of the subset, so by the multiplication rule
\begin{align*}
\#\{\text{subsets}\} = \underbrace{2\times 2\times\cdots\times 2}_{n \text{ times}} = 2^n
\end{align*}
\end{itemize}
Both expressions count the same set of subsets, so
\begin{align*}
\sum_{k=0}^n \binom{n}{k} = 2^n.
\end{align*}

%============================================================
\subsection{Book Ch.1, Exercise 16}
\begin{qbox}
Show that for all positive integers $n$ and $k$ with $n\ge k$,
\[
\binom{n}{k}+\binom{n}{k-1}=\binom{n+1}{k},
\]
doing this in two ways: (a) algebraically and (b) with a story, giving an interpretation for why both sides count the same thing.

\emph{Hint for the story proof:} Imagine $n+1$ people, with one of them pre-designated as ``president''.
\end{qbox}
\Sol
\textbf{(a) Algebraic proof.}
\begin{align*}
\binom{n}{k}+\binom{n}{k-1}
&= \frac{n!}{k!(n-k)!}+\frac{n!}{(k-1)!(n-k+1)!}\\[2pt]
&= \frac{n!(n-k+1)}{k!(n-k+1)!}+\frac{n!\,k}{k!(n-k+1)!} \qquad \text{(common denom.\ } k!(n-k+1)!\text{)}\\[2pt]
&= \frac{n!\big[(n-k+1)+k\big]}{k!(n-k+1)!}\\[2pt]
&= \frac{n!(n+1)}{k!(n-k+1)!}\\[2pt]
&= \frac{(n+1)!}{k!\big[(n+1)-k\big]!}\\[2pt]
&= \binom{n+1}{k}
\end{align*}
\textbf{(b) Story proof.} Consider $n+1$ people, one designated ``president''. The right side $\binom{n+1}{k}$ counts all size-$k$ subsets of these $n+1$ people. Split by whether the president is in the subset:
\begin{align*}
\underbrace{\binom{n}{k-1}}_{\text{president included: choose }k-1\text{ more from remaining }n} + \underbrace{\binom{n}{k}}_{\text{president excluded: choose all }k\text{ from remaining }n} = \binom{n+1}{k}
\end{align*}

%============================================================
\subsection{Book Ch.1, Exercise 18}
\begin{qbox}
(a) Show using a story proof that
\[
\binom{k}{k}+\binom{k+1}{k}+\binom{k+2}{k}+\cdots+\binom{n}{k}=\binom{n+1}{k+1},
\]
where $n$ and $k$ are positive integers with $n\ge k$. This is called the \emph{hockey stick identity}.

\emph{Hint:} Imagine arranging a group of people by age, and then think about the oldest person in a chosen subgroup.

(b) Suppose that a large pack of Haribo gummi bears can have anywhere between 30 and 50 gummi bears. There are 5 delicious flavors: pineapple (clear), raspberry (red), orange (orange), strawberry (green, mysteriously), and lemon (yellow). There are 0 non-delicious flavors. How many possibilities are there for the composition of such a pack of gummi bears? You can leave your answer in terms of a couple binomial coefficients, but not a sum of lots of binomial coefficients.
\end{qbox}
\Sol
\textbf{(a)} Consider choosing $k+1$ people out of $n+1$ people, and call the oldest person in the chosen subgroup ``Aemon''. Suppose there are $j$ people in the \emph{full} group younger than Aemon ($j$ ranges from $k$ to $n$, since Aemon needs at least $k$ younger people to fill out the rest of the subgroup). Then the other $k$ members of the subgroup must be chosen from those $j$ younger people:
\begin{align*}
\#\{\text{ways with Aemon as oldest, }j\text{ younger people below him}\} = \binom{j}{k}
\end{align*}
Summing over all valid positions for Aemon ($j=k,\dots,n$) counts every size-$(k+1)$ subgroup exactly once, which also equals $\binom{n+1}{k+1}$ directly:
\begin{align*}
\sum_{j=k}^{n}\binom{j}{k} = \binom{n+1}{k+1}
\end{align*}

\textbf{(b)} For a fixed pack size $i$, distributing $i$ (indistinguishable) bears among $5$ flavors is a stars-and-bars problem:
\begin{align*}
\#\{\text{compositions of }i\text{ bears}\} = \binom{5+i-1}{i} = \binom{i+4}{4}
\end{align*}
Summing over $i=30$ to $50$, and substituting $j=i+4$ (so $j$ runs from $34$ to $54$):
\begin{align*}
\sum_{i=30}^{50}\binom{i+4}{4} = \sum_{j=34}^{54}\binom{j}{4}
\end{align*}
Split the sum using the hockey-stick identity from part (a) with $k=4$:
\begin{align*}
\sum_{j=34}^{54}\binom{j}{4} &= \sum_{j=4}^{54}\binom{j}{4} - \sum_{j=4}^{33}\binom{j}{4}\\
&= \binom{55}{5} - \binom{34}{5}
\end{align*}
Now compute each term exactly:
\begin{align*}
\binom{55}{5} &= \frac{55\cdot 54\cdot 53\cdot 52\cdot 51}{5!} = \frac{417451320}{120} = 3478761\\
\binom{34}{5} &= \frac{34\cdot 33\cdot 32\cdot 31\cdot 30}{5!} = \frac{33390720}{120} = 278256
\end{align*}
\begin{align*}
\binom{55}{5}-\binom{34}{5} = 3478761 - 278256 = 3200505
\end{align*}
So there are $\boxed{3{,}200{,}505}$ possibilities.

%============================================================
\subsection{Book Ch.1, Exercise 22}
\begin{qbox}
A certain family has 6 children, consisting of 3 boys and 3 girls. Assuming that all birth orders are equally likely, what is the probability that the 3 eldest children are the 3 girls?
\end{qbox}
\Sol
Label the girls $1,2,3$ and boys $4,5,6$. A birth order is a permutation of $\{1,\dots,6\}$; there are $6!$ equally likely birth orders. The favorable outcomes are those where $\{1,2,3\}$ (in any order) come before $\{4,5,6\}$ (in any order):
\begin{align*}
\#\{\text{favorable orders}\} = 3! \times 3! = 6\times 6 = 36
\end{align*}
\begin{align*}
6! &= 720\\
P(\text{3 eldest are girls}) &= \frac{(3!)^2}{6!} = \frac{36}{720} = 0.05
\end{align*}
\emph{Check:} equivalently, there are $\binom{6}{3}=20$ equally likely ways to choose \emph{which} 3 birth-order slots the girls occupy, and only $1$ of these has the girls in the earliest $3$ slots, so $P=1/20=0.05$. \checkmark

%============================================================
\subsection{Book Ch.1, Exercise 23}
\begin{qbox}
A city with 6 districts has 6 robberies in a particular week. Assume the robberies are located randomly, with all possibilities for which robbery occurred where equally likely. What is the probability that some district had more than 1 robbery?
\end{qbox}
\Sol
Each of the 6 robberies independently lands in one of 6 districts, so there are $6^6$ equally likely configurations. The complement (``no district has more than 1'') means all 6 robberies land in different districts, i.e.\ a bijection between robberies and districts: $6!$ ways.
\begin{align*}
6^6 &= 46656\\
6! &= 720\\
P(\text{no repeats}) &= \frac{6!}{6^6} = \frac{720}{46656} = 0.0154321\ldots
\end{align*}
\begin{align*}
P(\text{some district}>1\text{ robbery}) = 1 - \frac{6!}{6^6} = 1-0.0154321 \approx 0.9846
\end{align*}

%============================================================
\subsection{Book Ch.1, Exercise 26}
\begin{qbox}
A college has 10 (non-overlapping) time slots for its courses, and blithely assigns courses to time slots randomly and independently. A student randomly chooses 3 of the courses to enroll in. What is the probability that there is a conflict in the student's schedule?
\end{qbox}
\Sol
For no conflict, the 3 chosen courses must land in 3 different (of 10) time slots. Assign slots to the 3 courses one at a time:
\begin{align*}
P(\text{no conflict}) = \frac{10}{10}\cdot\frac{9}{10}\cdot\frac{8}{10} = \frac{10\times 9\times 8}{10^3} = \frac{720}{1000} = 0.72
\end{align*}
\begin{align*}
P(\text{conflict}) = 1-0.72 = 0.28
\end{align*}

%============================================================
\subsection{Book Ch.1, Exercise 27}
\begin{qbox}
For each part, decide whether the blank should be filled in with $=$, $<$, or $>$, and give a clear explanation.

(a) (probability that the total after rolling 4 fair dice is 21) \underline{\hspace{1cm}} (probability that the total after rolling 4 fair dice is 22)

(b) (probability that a random 2-letter word is a palindrome) \underline{\hspace{1cm}} (probability that a random 3-letter word is a palindrome)
\end{qbox}
\Sol
\textbf{(a)} All $6^4$ \emph{ordered} outcomes of 4 dice are equally likely, so we count ordered outcomes summing to each target.

Sum $=21$ comes from unordered multisets $(6,6,6,3)$, $(6,5,5,5)$, $(6,6,5,4)$:
\begin{align*}
\#(6,6,6,3) &= \frac{4!}{3!\,1!} = 4\\
\#(6,5,5,5) &= \frac{4!}{1!\,3!} = 4\\
\#(6,6,5,4) &= \frac{4!}{2!\,1!\,1!} = 12\\
\text{Total for }21 &= 4+4+12 = 20
\end{align*}
Sum $=22$ comes from $(6,6,6,4)$, $(6,6,5,5)$:
\begin{align*}
\#(6,6,6,4) &= \frac{4!}{3!\,1!} = 4\\
\#(6,6,5,5) &= \frac{4!}{2!\,2!} = 6\\
\text{Total for }22 &= 4+6 = 10
\end{align*}
Since $20>10$ (in fact exactly double),
\begin{align*}
P(\text{sum}=21) \;>\; P(\text{sum}=22)
\end{align*}

\textbf{(b)} For a random word of length $L$ over a 26-letter alphabet, being a palindrome means the first $\lceil L/2\rceil$ letters are free and the rest are determined by symmetry:
\begin{align*}
P(\text{palindrome, length }L) = \frac{26^{\lceil L/2\rceil}}{26^{L}} = 26^{\lceil L/2\rceil - L}
\end{align*}
For $L=2$: $\lceil 2/2\rceil - 2 = 1-2=-1$, so $P=26^{-1}$.
For $L=3$: $\lceil 3/2\rceil - 3 = 2-3=-1$, so $P=26^{-1}$.
\begin{align*}
P(\text{2-letter palindrome}) = P(\text{3-letter palindrome}) = \frac{1}{26}
\end{align*}
so the blank is $=$.

%============================================================
\subsection{Book Ch.1, Exercise 29}
\begin{qbox}
Elk dwell in a certain forest. There are $N$ elk, of which a simple random sample of size $n$ are captured and tagged (``simple random sample'' means that all $\binom{N}{n}$ sets of $n$ elk are equally likely). The captured elk are returned to the population, and then a new sample is drawn, this time with size $m$. This is an important method that is widely used in ecology, known as \emph{capture-recapture}. What is the probability that exactly $k$ of the $m$ elk in the new sample were previously tagged? (Assume that an elk that was captured before doesn't become more or less likely to be captured again.)
\end{qbox}
\Sol
All $\binom{N}{m}$ samples of size $m$ are equally likely. To get exactly $k$ tagged elk in the new sample, choose $k$ of the $n$ tagged elk \emph{and} $m-k$ of the $N-n$ untagged elk:
\begin{align*}
\#\{\text{favorable samples}\} = \binom{n}{k}\binom{N-n}{m-k}
\end{align*}
\begin{align*}
P(\text{exactly }k\text{ tagged}) = \frac{\dbinom{n}{k}\dbinom{N-n}{m-k}}{\dbinom{N}{m}}
\end{align*}
valid for $0\le k\le n$ and $0\le m-k\le N-n$ (else the probability is $0$). This is the \emph{Hypergeometric} distribution.

%============================================================
\subsection{Book Ch.1, Exercise 31}
\begin{qbox}
A jar contains $r$ red balls and $g$ green balls, where $r$ and $g$ are fixed positive integers. A ball is drawn from the jar randomly (with all possibilities equally likely), and then a second ball is drawn randomly.

(a) Explain intuitively why the probability of the second ball being green is the same as the probability of the first ball being green.

(b) Define notation for the sample space of the problem, and use this to compute the probabilities from (a) and show that they are the same.

(c) Suppose that there are 16 balls in total, and that the probability that the two balls are the same color is the same as the probability that they are different colors. What are $r$ and $g$ (list all possibilities)?
\end{qbox}
\Sol
\textbf{(a)} By symmetry, none of the $g+r$ balls is special for the ``second draw'' — we could equally well imagine drawing both balls simultaneously with two hands and then arbitrarily calling one ``first'' and the other ``second''. So each ball is equally likely to occupy either position.

\textbf{(b)} Label balls $1,\dots,g+r$, with $1,\dots,g$ green. Sample space $=\{(a,b): a,b\in\{1,\dots,g+r\},\ a\ne b\}$, each pair equally likely. Let $G_i=\{i\text{th ball drawn is green}\}$.
\begin{align*}
\#\{\text{sample space}\} &= (g+r)(g+r-1)\\
\#\{G_1\} &= g\cdot(g+r-1) \quad\text{(green first, any of remaining $g+r-1$ second)}\\
\#\{G_2\} &= (g+r-1)\cdot g \quad\text{(any first, green second: $g$ favorable seconds $\times(g+r-1)$ firsts)}
\end{align*}
\begin{align*}
P(G_1)=P(G_2) = \frac{g(g+r-1)}{(g+r)(g+r-1)} = \frac{g}{g+r}
\end{align*}

\textbf{(c)} Let $A=\{$one ball of each color$\}=(G_1\cap G_2^c)\cup(G_1^c\cap G_2)$. Number of favorable ordered pairs:
\begin{align*}
\#A = \underbrace{g\cdot r}_{\text{green then red}} + \underbrace{r\cdot g}_{\text{red then green}} = 2gr
\end{align*}
\begin{align*}
P(A) = \frac{2gr}{(g+r)(g+r-1)}
\end{align*}
We're told $P(A)=P(A^c)$, so $P(A)=\tfrac12$. With $g+r=16$:
\begin{align*}
\frac{2gr}{16\times 15} &= \frac12\\
\frac{2gr}{240} &= \frac12\\
2gr &= 120\\
gr &= 60
\end{align*}
So $g,r$ are roots of $t^2-16t+60=0$:
\begin{align*}
t = \frac{16\pm\sqrt{16^2-4(60)}}{2} = \frac{16\pm\sqrt{256-240}}{2} = \frac{16\pm\sqrt{16}}{2} = \frac{16\pm4}{2}
\end{align*}
\begin{align*}
t = 10 \quad\text{or}\quad t=6
\end{align*}
So $(g,r)=(10,6)$ or $(g,r)=(6,10)$.

%============================================================
\subsection{Book Ch.1, Exercise 32}
\begin{qbox}
A random 5-card poker hand is dealt from a standard deck of cards. Find the probability of each of the following possibilities (in terms of binomial coefficients).

(a) A flush (all 5 cards being of the same suit; do not count a royal flush, which is a flush with an ace, king, queen, jack, and 10).

(b) Two pair (e.g., two 3's, two 7's, and an ace).
\end{qbox}
\Sol
Total hands: $\binom{52}{5}=2598960$.

\textbf{(a)} Choose the suit ($4$ ways), then $5$ of that suit's $13$ cards, minus the $1$ way to get the royal flush in that suit:
\begin{align*}
\binom{13}{5} &= 1287\\
\binom{13}{5}-1 &= 1286\\
4\left(\binom{13}{5}-1\right) &= 4\times 1286 = 5144
\end{align*}
\begin{align*}
P(\text{flush}) = \frac{4\left[\binom{13}{5}-1\right]}{\binom{52}{5}} = \frac{5144}{2598960} \approx 0.001979
\end{align*}

\textbf{(b)} Choose 2 ranks for the pairs $\binom{13}{2}$, choose 2 suits out of 4 for each paired rank $\binom{4}{2}$ each, then choose 1 of the remaining $52-8=44$ cards for the kicker:
\begin{align*}
\binom{13}{2} &= 78\\
\binom{4}{2}^2 &= 6^2 = 36\\
78\times 36 &= 2808\\
2808\times 44 &= 123552
\end{align*}
\begin{align*}
P(\text{two pair}) = \frac{\binom{13}{2}\binom{4}{2}^2\cdot 44}{\binom{52}{5}} = \frac{123552}{2598960} \approx 0.047539
\end{align*}

%============================================================
\subsection{Book Ch.1, Exercise 40}
\begin{qbox}
A \emph{norepeatword} is a sequence of at least one (and possibly all) of the usual 26 letters a,b,c,\dots,z, with repetitions not allowed. For example, ``course'' is a norepeatword, but ``statistics'' is not. Order matters, e.g., ``course'' is not the same as ``source''.

A norepeatword is chosen randomly, with all norepeatwords equally likely. Show that the probability that it uses all 26 letters is very close to $1/e$.
\end{qbox}
\Sol
Norepeatwords with all 26 letters: ordered arrangements of all 26 letters, i.e.\ $26!$.

Norepeatwords with exactly $k$ letters ($1\le k\le 26$): choose which $k$ letters ($\binom{26}{k}$ ways), then arrange them ($k!$ ways):
\begin{align*}
\#\{\text{length-}k\text{ norepeatwords}\} = \binom{26}{k}k!
\end{align*}
So the total number of norepeatwords is $\displaystyle\sum_{k=1}^{26}\binom{26}{k}k!$, and
\begin{align*}
P(\text{all 26 letters}) &= \frac{26!}{\displaystyle\sum_{k=1}^{26}\binom{26}{k}k!}\\
&= \frac{26!}{\displaystyle\sum_{k=1}^{26}\frac{26!}{k!(26-k)!}k!}\\
&= \frac{26!}{\displaystyle\sum_{k=1}^{26}\frac{26!}{(26-k)!}} \qquad \text{(the $k!$ cancels)}
\end{align*}
Divide numerator and denominator by $26!$, and substitute $m=26-k$ (so $m$ runs from $0$ to $25$ as $k$ runs from $26$ down to $1$):
\begin{align*}
P(\text{all 26 letters}) = \frac{1}{\displaystyle\sum_{k=1}^{26}\frac{1}{(26-k)!}} = \frac{1}{\dfrac{1}{25!}+\dfrac{1}{24!}+\cdots+\dfrac{1}{1!}+\dfrac{1}{0!}}
\end{align*}
The denominator is exactly the first 26 terms ($m=0$ to $25$) of the Taylor series
\begin{align*}
e^x = \sum_{m=0}^\infty \frac{x^m}{m!}\Big|_{x=1} = \sum_{m=0}^\infty \frac{1}{m!}
\end{align*}
Since the series for $e$ converges extremely fast (the tail past $m=25$ is smaller than $10^{-26}$), the denominator $\approx e$, so
\begin{align*}
P(\text{all 26 letters}) \approx \frac{1}{e} \approx 0.367879
\end{align*}

%============================================================
\subsection{Book Ch.1, Exercise 46}
\begin{qbox}
Arby has a belief system assigning a number $P_{\text{Arby}}(A)$ between 0 and 1 to every event $A$ (for some sample space). This represents Arby's degree of belief about how likely $A$ is to occur. For any event $A$, Arby is willing to pay a price of $1000\cdot P_{\text{Arby}}(A)$ dollars to buy a certificate such as the one shown below:

\medskip
\fbox{\parbox{0.85\linewidth}{\textbf{Certificate}\\[2pt]
The owner of this certificate can redeem it for \$1000 if $A$ occurs. No value if $A$ does not occur, except as required by federal, state, or local law. No expiration date.}}
\medskip

Likewise, Arby is willing to sell such a certificate at the same price. Indeed, Arby is willing to buy or sell any number of certificates at this price, as Arby considers it the ``fair'' price.

Arby stubbornly refuses to accept the axioms of probability. In particular, suppose that there are two disjoint events $A$ and $B$ with
\[
P_{\text{Arby}}(A\cup B) \ne P_{\text{Arby}}(A) + P_{\text{Arby}}(B).
\]
Show how to make Arby go bankrupt, by giving a list of transactions Arby is willing to make that will guarantee that Arby will lose money (you can assume it will be known whether $A$ occurred and whether $B$ occurred the day after any certificates are bought/sold).
\end{qbox}
\Sol
Work in units of thousands of dollars, so a certificate on event $C$ costs $P_{\text{Arby}}(C)$.

\textbf{Case 1: $P_{\text{Arby}}(A\cup B) < P_{\text{Arby}}(A)+P_{\text{Arby}}(B)$.}
Buy one $A$-certificate and one $B$-certificate; sell one $(A\cup B)$-certificate.
\begin{align*}
\text{Immediate net cost to Arby's counterparty} &= P_{\text{Arby}}(A)+P_{\text{Arby}}(B) - P_{\text{Arby}}(A\cup B) > 0
\end{align*}
Since $A,B$ are disjoint, exactly one of the following happens:
\begin{itemize}[leftmargin=1.3em]
\item Neither $A$ nor $B$ occurs: all three certificates pay \$0 — Arby's initial loss stands.
\item Exactly one of $A,B$ occurs (say $A$): Arby receives \$1000 on the $A$-certificate he sold... wait, he \emph{bought} the $A$-certificate, so he receives \$1000; but $A\cup B$ also occurred, so Arby (who \emph{sold} the $(A\cup B)$-certificate) must pay out \$1000. These cancel exactly.
\end{itemize}
So Arby's payouts always net to zero, and Arby is left having paid the strictly positive amount $P_{\text{Arby}}(A)+P_{\text{Arby}}(B)-P_{\text{Arby}}(A\cup B)$ with certainty.

\textbf{Case 2: $P_{\text{Arby}}(A\cup B) > P_{\text{Arby}}(A)+P_{\text{Arby}}(B)$.}
Reverse the trades: sell an $A$-certificate, sell a $B$-certificate, buy an $(A\cup B)$-certificate.
\begin{align*}
\text{Immediate net gain to Arby's counterparty} &= P_{\text{Arby}}(A\cup B) - \big[P_{\text{Arby}}(A)+P_{\text{Arby}}(B)\big] > 0
\end{align*}
Again, since $A,B$ disjoint, any payout Arby owes on the sold certificates is exactly matched by what Arby receives on the $(A\cup B)$-certificate he bought. Arby loses the gap with certainty.

Repeating these trades $n$ times multiplies the guaranteed loss by $n$, so Arby can be driven to bankruptcy — an \emph{arbitrage opportunity}. This shows the additivity axiom is not arbitrary but necessary for coherent belief.

%============================================================
\subsection{Book Ch.1, Exercise 48}
\begin{qbox}
A card player is dealt a 13-card hand from a well-shuffled, standard deck of cards. What is the probability that the hand is void in at least one suit (``void in a suit'' means having no cards of that suit)?
\end{qbox}
\Sol
Let $S,H,D,C$ be the events of being void in Spades, Hearts, Diamonds, Clubs respectively. By inclusion--exclusion and symmetry among the 4 suits:
\begin{align*}
P(S\cup H\cup D\cup C) = 4P(S) - 6P(S\cap H) + 4P(S\cap H\cap D) - P(S\cap H\cap D\cap C)
\end{align*}
(coefficients $\binom{4}{1}=4$, $\binom{4}{2}=6$, $\binom{4}{3}=4$, $\binom{4}{4}=1$). The last term is $0$ (can't be void in all 4 suits with 13 cards). Void in one specific suit means all 13 cards come from the other $39$; void in two specific suits means all 13 come from the other $26$; void in three specific suits means all 13 come from the remaining $13$:
\begin{align*}
P(S) &= \frac{\binom{39}{13}}{\binom{52}{13}}, \qquad
P(S\cap H) = \frac{\binom{26}{13}}{\binom{52}{13}}, \qquad
P(S\cap H\cap D) = \frac{\binom{13}{13}}{\binom{52}{13}} = \frac{1}{\binom{52}{13}}
\end{align*}
\begin{align*}
\binom{39}{13} &= 8122425444, \qquad \binom{26}{13} = 10400600, \qquad \binom{52}{13} = 635013559600
\end{align*}
\begin{align*}
P(S\cup H\cup D\cup C) &= \frac{4(8122425444) - 6(10400600) + 4(1)}{635013559600}\\
&= \frac{32489701776 - 62403600 + 4}{635013559600}\\
&= \frac{32427298180}{635013559600}\\
&\approx 0.05107
\end{align*}

%============================================================
\subsection{Book Ch.1, Exercise 52}
\begin{qbox}
Alice attends a small college in which each class meets only once a week. She is deciding between 30 non-overlapping classes. There are 6 classes to choose from for each day of the week, Monday through Friday. Trusting in the benevolence of randomness, Alice decides to register for 7 randomly selected classes out of the 30, with all choices equally likely. What is the probability that she will have classes every day, Monday through Friday? (This problem can be done either directly using the naive definition of probability, or using inclusion-exclusion.)
\end{qbox}
\Sol
Total ways to pick 7 of 30 classes: $\binom{30}{7}=2035800$.

\textbf{Direct counting.} With 7 classes over 5 days, each day $\ge1$ class, the only possible day-distributions of class-counts are: (i) two days with 2 classes, three days with 1 class; or (ii) one day with 3 classes, four days with 1 class.

\emph{Case (i):} choose which 2 of the 5 days get 2 classes $\binom{5}{2}$; for each of those 2 days choose 2 of its 6 classes $\binom{6}{2}$ each; for each of the other 3 days choose 1 of 6 classes ($6$ each):
\begin{align*}
\binom{5}{2}\binom{6}{2}^2 6^3 = 10\times 15^2\times 216 = 10\times225\times216 = 486000
\end{align*}
\emph{Case (ii):} choose which day gets 3 classes $\binom{5}{1}$, choose 3 of its 6 classes $\binom{6}{3}$, and 1 class from each of the other 4 days ($6^4$):
\begin{align*}
\binom{5}{1}\binom{6}{3}6^4 = 5\times 20\times 1296 = 129600
\end{align*}
\begin{align*}
P(\text{class every day}) = \frac{486000+129600}{\binom{30}{7}} = \frac{615600}{2035800} = \frac{114}{377} \approx 0.302
\end{align*}

\textbf{Check via inclusion-exclusion on the complement.} Let $B_i=\{$no class on day $i\}$; we want $1-P(B_1\cup\cdots\cup B_5)$. If Alice avoids $j$ specific days entirely, she must pick all 7 classes from the remaining $30-6j$ classes:
\begin{align*}
P(B_1) = \frac{\binom{24}{7}}{\binom{30}{7}},\qquad P(B_1\cap B_2) = \frac{\binom{18}{7}}{\binom{30}{7}}, \qquad P(B_1\cap B_2\cap B_3) = \frac{\binom{12}{7}}{\binom{30}{7}}
\end{align*}
(terms with $4$ or more $B_i$'s vanish, since that would leave only $6$ classes for $7$ picks). With $\binom{24}{7}=346104$, $\binom{18}{7}=31824$, $\binom{12}{7}=792$:
\begin{align*}
P(B_1\cup\cdots\cup B_5) &= 5\binom{24}{7}\Big/\binom{30}{7} - \binom{5}{2}\binom{18}{7}\Big/\binom{30}{7} + \binom{5}{3}\binom{12}{7}\Big/\binom{30}{7}\\
&= \frac{5(346104) - 10(31824) + 10(792)}{2035800}\\
&= \frac{1730520 - 318240 + 7920}{2035800} = \frac{1420200}{2035800} = \frac{263}{377}
\end{align*}
\begin{align*}
P(\text{class every day}) = 1-\frac{263}{377} = \frac{114}{377} \approx 0.302 \quad\checkmark
\end{align*}

%============================================================
\subsection{Book Ch.1, Exercise 59}
\begin{qbox}
There are 100 passengers lined up to board an airplane with 100 seats (with each seat assigned to one of the passengers). The first passenger in line crazily decides to sit in a randomly chosen seat (with all seats equally likely). Each subsequent passenger takes his or her assigned seat if available, and otherwise sits in a random available seat. What is the probability that the last passenger in line gets to sit in his or her assigned seat? (This is a common interview problem, and a beautiful example of the power of symmetry.)

\emph{Hint:} Call the seat assigned to the $j$th passenger in line ``seat $j$'' (regardless of whether the airline calls it seat 23A or whatever). What are the possibilities for which seats are available to the last passenger in line, and what is the probability of each of these possibilities?
\end{qbox}
\Sol
Claim: the only seats that can possibly be available for the last (100th) passenger are seat 1 and seat 100. Indeed, for any seat $j$ with $2\le j\le 99$: passenger $j$ would have sat in seat $j$ if it were available at their turn, and if it wasn't, someone before them took it — either way, seat $j$ gets taken before the last passenger's turn (either by passenger $j$ or by whoever displaced them). So seat $j$ is never what's left over.

By symmetry, among the first 99 passengers, seat $1$ and seat $100$ are interchangeable in terms of who might end up displacing whom — neither seat is "special" until someone actually chooses one of them. So
\begin{align*}
P(\text{seat 1 available at the end}) = P(\text{seat 100 available at the end})
\end{align*}
and since exactly one of these two must be available (and they're the only options):
\begin{align*}
P(\text{seat 1 avail.}) + P(\text{seat 100 avail.}) = 1 \implies P(\text{seat 100 avail.}) = \frac12
\end{align*}
So the probability the last passenger gets their own assigned seat (seat 100) is
\begin{align*}
P = \frac12
\end{align*}

\newpage
\section*{Strategic Practice}
\addcontentsline{toc}{section}{Strategic Practice}

%============================================================
\subsection{SP 1 §1 · Problem 1}
\begin{qbox}
For each part, decide whether the blank should be filled in with $=$, $<$, or $>$, and give a short but clear explanation.

(a) (probability that the total after rolling 4 fair dice is 21) \underline{\hspace{1cm}} (probability that the total after rolling 4 fair dice is 22)

(b) (probability that a random 2 letter word is a palindrome) \underline{\hspace{1cm}} (probability that a random 3 letter word is a palindrome)
\end{qbox}
\Sol
Identical to Q9. Recall the ordered-outcome counts:
\begin{align*}
\#\{\text{sum}=21\} &= \underbrace{4}_{(6,6,6,3)}+\underbrace{4}_{(6,5,5,5)}+\underbrace{12}_{(6,6,5,4)} = 20\\
\#\{\text{sum}=22\} &= \underbrace{4}_{(6,6,6,4)}+\underbrace{6}_{(6,6,5,5)} = 10
\end{align*}
\begin{align*}
\text{(a)}\ P(21) > P(22) \qquad\qquad \text{(b)}\ P(\text{2-letter pal.}) = P(\text{3-letter pal.}) = \frac{1}{26}
\end{align*}

%============================================================
\subsection{SP 1 §1 · Problem 2}
\begin{qbox}
A random 5 card poker hand is dealt from a standard deck of cards. Find the probability of each of the following (in terms of binomial coefficients).

(a) A flush (all 5 cards being of the same suit; do not count a royal flush, which is a flush with an Ace, King, Queen, Jack, and 10)

(b) Two pair (e.g., two 3's, two 7's, and an Ace)
\end{qbox}
\Sol
Identical setup to Q12.
\begin{align*}
P(\text{flush}) &= \frac{4\left[\binom{13}{5}-1\right]}{\binom{52}{5}} = \frac{4(1287-1)}{2598960}=\frac{5144}{2598960}\approx 0.001979\\[4pt]
P(\text{two pair}) &= \frac{\binom{13}{2}\binom{4}{2}^2\cdot 44}{\binom{52}{5}} = \frac{78\times 36\times 44}{2598960}=\frac{123552}{2598960}\approx 0.047539
\end{align*}

%============================================================
\subsection{SP 1 §1 · Problem 3}
\begin{qbox}
(a) How many paths are there from the point $(0, 0)$ to the point $(110, 111)$ in the plane such that each step either consists of going one unit up or one unit to the right?

(b) How many paths are there from $(0, 0)$ to $(210, 211)$, where each step consists of going one unit up or one unit to the right, and the path has to go through $(110, 111)$?
\end{qbox}
\Sol
Identical to Q2.
\begin{align*}
\text{(a)}\quad \binom{221}{110} \qquad\qquad
\text{(b)}\quad \binom{221}{110}\binom{200}{100}
\end{align*}

%============================================================
\subsection{SP 1 §1 · Problem 4}
\begin{qbox}
A norepeatword is a sequence of at least one (and possibly all) of the usual 26 letters a,b,c,\dots,z, with repetitions not allowed. For example, ``course'' is a norepeatword, but ``statistics'' is not. Order matters, e.g., ``course'' is not the same as ``source''.

A norepeatword is chosen randomly, with all norepeatwords equally likely. Show that the probability that it uses all 26 letters is very close to $1/e$.
\end{qbox}
\Sol
Identical to Q13.
\begin{align*}
P(\text{all 26 letters}) = \frac{1}{\dfrac{1}{25!}+\dfrac{1}{24!}+\cdots+\dfrac{1}{1!}+\dfrac{1}{0!}} \approx \frac{1}{e} \approx 0.367879
\end{align*}

%============================================================
\subsection{SP 1 §2 · Problem 5}
\begin{qbox}
Give a story proof that $\displaystyle\sum_{k=0}^n \binom{n}{k}=2^n$.
\end{qbox}
\Sol
Identical to Q3.
\begin{align*}
\underbrace{\sum_{k=0}^n \binom{n}{k}}_{\text{subsets counted by size}} = \underbrace{2^n}_{\text{each of $n$ people in/out independently}}
\end{align*}

%============================================================
\subsection{SP 1 §2 · Problem 6}
\begin{qbox}
Give a story proof that
\[
\frac{(2n)!}{2^n \cdot n!} = (2n-1)(2n-3)\cdots 3\cdot 1.
\]
\end{qbox}
\Sol
Consider forming $n$ (unordered) partnerships out of $2n$ people.

\textbf{Method 1.} Line up all $2n$ people ($(2n)!$ orderings); declare the 1st\&2nd a pair, 3rd\&4th a pair, etc. This overcounts each partnership-set by $2^n$ (order within each of the $n$ pairs doesn't matter) times $n!$ (order among the $n$ pairs doesn't matter):
\begin{align*}
\#\{\text{partnership-sets}\} = \frac{(2n)!}{2^n\cdot n!}
\end{align*}

\textbf{Method 2 (direct).} Person 1 has $2n-1$ possible partners. Once that pair is fixed, take the next unpaired person: they have $2n-3$ possible partners (everyone except themselves and the already-paired 2 people). Continuing, the next unpaired person has $2n-5$ choices, and so on down to the last pair, which has $1$ choice:
\begin{align*}
\#\{\text{partnership-sets}\} = (2n-1)(2n-3)(2n-5)\cdots 3\cdot 1
\end{align*}

Both methods count the same set of partnership-configurations, so
\begin{align*}
\frac{(2n)!}{2^n\cdot n!} = (2n-1)(2n-3)\cdots 3\cdot 1
\end{align*}

%============================================================
\subsection{SP 1 §2 · Problem 7}
\begin{qbox}
Show that for all positive integers $n$ and $k$ with $n\ge k$,
\[
\binom{n}{k}+\binom{n}{k-1}=\binom{n+1}{k},
\]
doing this in two ways: (a) algebraically and (b) with a ``story'', giving an interpretation for why both sides count the same thing.

\emph{Hint for the ``story'' proof:} imagine $n+1$ people, with one of them pre-designated as ``president''.
\end{qbox}
\Sol
Identical to Q4.
\begin{align*}
\binom{n}{k}+\binom{n}{k-1} &= \frac{n!}{k!(n-k)!}+\frac{n!}{(k-1)!(n-k+1)!}
= \frac{n!(n-k+1)+n!\,k}{k!(n-k+1)!} = \frac{n!(n+1)}{k!(n-k+1)!} = \binom{n+1}{k}
\end{align*}
Story: split size-$k$ subgroups of $n+1$ people (with a designated president) by whether the president is in $(\binom{n}{k-1}$ ways$)$ or out $(\binom{n}{k}$ ways$)$.

\newpage
\subsection{SP 2 §1 · Problem 1}
\begin{qbox}
For a group of 7 people, find the probability that all 4 seasons (winter, spring, summer, fall) occur at least once each among their birthdays, assuming that all seasons are equally likely.
\end{qbox}
\Sol
Let $A_i=\{$no birthday falls in season $i\}$, $i=1,\dots,4$. We want $1-P(A_1\cup A_2\cup A_3\cup A_4)$. By inclusion--exclusion and symmetry among seasons (note $A_1\cap A_2\cap A_3\cap A_4=\emptyset$, since someone must be born in \emph{some} season):
\begin{align*}
P(A_1\cup A_2\cup A_3\cup A_4) = 4P(A_1) - 6P(A_1\cap A_2) + 4P(A_1\cap A_2\cap A_3)
\end{align*}
Each person avoids a single forbidden season with probability $3/4$; avoids 2 forbidden seasons with probability $2/4=1/2$; avoids 3 forbidden seasons with probability $1/4$ — each independently across the 7 people:
\begin{align*}
P(A_1) = \left(\frac34\right)^7, \qquad P(A_1\cap A_2) = \left(\frac12\right)^7 = \frac{1}{2^7}, \qquad P(A_1\cap A_2\cap A_3) = \left(\frac14\right)^7 = \frac{1}{4^7}
\end{align*}
Numerically:
\begin{align*}
4\left(\frac34\right)^7 &= 4(0.1334839\ldots) = 0.5339355\\
\frac{6}{2^7} &= \frac{6}{128} = 0.0468750\\
\frac{4}{4^7} &= \frac{4}{16384} = 0.0002441
\end{align*}
\begin{align*}
P(A_1\cup A_2\cup A_3\cup A_4) = 0.5339355 - 0.0468750 + 0.0002441 = 0.4873047
\end{align*}
\begin{align*}
P(\text{all 4 seasons occur}) = 1-0.4873047 \approx 0.5127
\end{align*}

%============================================================
\subsection{SP 2 §1 · Problem 2}
\begin{qbox}
Alice attends a small college in which each class meets only once a week. She is deciding between 30 non-overlapping classes. There are 6 classes to choose from for each day of the week, Monday through Friday. Trusting in the benevolence of randomness, Alice decides to register for 7 randomly selected classes out of the 30, with all choices equally likely. What is the probability that she will have classes every day, Monday through Friday? (This problem can be done either directly using the naive definition of probability, or using inclusion-exclusion.)
\end{qbox}
\Sol
Identical to Q16.
\begin{align*}
P(\text{class every day}) = \frac{\binom{5}{2}\binom{6}{2}^26^3+\binom{5}{1}\binom{6}{3}6^4}{\binom{30}{7}} = \frac{486000+129600}{2035800} = \frac{615600}{2035800}=\frac{114}{377}\approx 0.302
\end{align*}

\newpage
\section*{Homework}
\addcontentsline{toc}{section}{Homework}

%============================================================
\subsection{HW 1 · Problem 1}
\begin{qbox}
A certain family has 6 children, consisting of 3 boys and 3 girls. Assuming that all birth orders are equally likely, what is the probability that the 3 eldest children are the 3 girls?
\end{qbox}
\Sol
Identical to Q6.
\begin{align*}
P = \frac{(3!)^2}{6!} = \frac{36}{720} = 0.05
\end{align*}

%============================================================
\subsection{HW 1 · Problem 2}
\begin{qbox}
(a) How many ways are there to split a dozen people into 3 teams, where one team has 2 people, and the other two teams have 5 people each?

(b) How many ways are there to split a dozen people into 3 teams, where each team has 4 people?
\end{qbox}
\Sol
Identical to Q1.
\begin{align*}
\text{(a)}\quad \frac{\binom{12}{2}\binom{10}{5}}{2} = \frac{66\times 252}{2} = \frac{16632}{2} = 8316
\end{align*}
\begin{align*}
\text{(b)}\quad \frac{12!}{4!\,4!\,4!\,3!} = \frac{479001600}{82944} = 5775
\end{align*}

%============================================================
\subsection{HW 1 · Problem 3}
\begin{qbox}
A college has 10 (non-overlapping) time slots for its courses, and blithely assigns courses to time slots randomly and independently. A student randomly chooses 3 of the courses to enroll in (for the PTP, to avoid getting fined). What is the probability that there is a conflict in the student's schedule?
\end{qbox}
\Sol
Identical to Q8.
\begin{align*}
P(\text{no conflict}) = \frac{10\times9\times8}{10^3} = \frac{720}{1000}=0.72, \qquad P(\text{conflict}) = 1-0.72 = 0.28
\end{align*}

%============================================================
\subsection{HW 1 · Problem 4}
\begin{qbox}
A city with 6 districts has 6 robberies in a particular week. Assume the robberies are located randomly, with all possibilities for which robbery occurred where equally likely. What is the probability that some district had more than 1 robbery?
\end{qbox}
\Sol
Identical to Q7.
\begin{align*}
P(\text{some district}>1) = 1-\frac{6!}{6^6} = 1-\frac{720}{46656} = 1-0.015432 \approx 0.9846
\end{align*}

%============================================================
\subsection{HW 1 · Problem 5}
\begin{qbox}
Elk dwell in a certain forest. There are $N$ elk, of which a simple random sample of size $n$ are captured and tagged (``simple random sample'' means that all $\binom{N}{n}$ sets of $n$ elk are equally likely). The captured elk are returned to the population, and then a new sample is drawn, this time with size $m$. This is an important method that is widely-used in ecology, known as \emph{capture-recapture}.

What is the probability that exactly $k$ of the $m$ elk in the new sample were previously tagged? (Assume that an elk that was captured before doesn't become more or less likely to be captured again.)
\end{qbox}
\Sol
Identical to Q10.
\begin{align*}
P(\text{exactly }k\text{ tagged}) = \frac{\dbinom{n}{k}\dbinom{N-n}{m-k}}{\dbinom{N}{m}}
\end{align*}

%============================================================
\subsection{HW 1 · Problem 6}
\begin{qbox}
A jar contains $r$ red balls and $g$ green balls, where $r$ and $g$ are fixed positive integers. A ball is drawn from the jar randomly (with all possibilities equally likely), and then a second ball is drawn randomly.

(a) Explain intuitively why the probability of the second ball being green is the same as the probability of the first ball being green.

(b) Define notation for the sample space of the problem, and use this to compute the probabilities from (a) and show that they are the same.

(c) Suppose that there are 16 balls in total, and that the probability that the two balls are the same color is the same as the probability that they are different colors. What are $r$ and $g$ (list all possibilities)?
\end{qbox}
\Sol
Identical to Q11.
\begin{align*}
P(G_1) = P(G_2) = \frac{g}{g+r}, \qquad gr = 60,\ g+r=16 \implies (g,r) = (10,6)\ \text{or}\ (6,10)
\end{align*}

%============================================================
\subsection{HW 1 · Problem 7}
\begin{qbox}
(a) Show using a story proof that
\[
\binom{k}{k}+\binom{k+1}{k}+\binom{k+2}{k}+\cdots+\binom{n}{k}=\binom{n+1}{k+1},
\]
where $n$ and $k$ are positive integers with $n\ge k$.

\emph{Hint:} imagine arranging a group of people by age, and then think about the oldest person in a chosen subgroup.

(b) Suppose that a large pack of Haribo gummi bears can have anywhere between 30 and 50 gummi bears. There are 5 delicious flavors: pineapple (clear), raspberry (red), orange (orange), strawberry (green, mysteriously), and lemon (yellow). There are 0 non-delicious flavors. How many possibilities are there for the composition of such a pack of gummi bears? You can leave your answer in terms of a couple binomial coefficients, but not a sum of lots of binomial coefficients.
\end{qbox}
\Sol
Identical to Q5.
\begin{align*}
\sum_{j=k}^n \binom{j}{k} &= \binom{n+1}{k+1}\\[4pt]
\sum_{i=30}^{50}\binom{i+4}{4} &= \binom{55}{5}-\binom{34}{5} = 3478761-278256 = 3200505
\end{align*}

\newpage
%============================================================
\subsection{HW 2 · Problem 1}
\begin{qbox}
Arby has a belief system assigning a number $P_{\text{Arby}}(A)$ between 0 and 1 to every event $A$ (for some sample space). This represents Arby's subjective degree of belief about how likely $A$ is to occur. For any event $A$, Arby is willing to pay a price of $1000\cdot P_{\text{Arby}}(A)$ dollars to buy a certificate such as the one shown below:

\medskip
\fbox{\parbox{0.85\linewidth}{\textbf{Certificate}\\[2pt]
The owner of this certificate can redeem it for \$1000 if $A$ occurs. No value if $A$ does not occur, except as required by federal, state, or local law. No expiration date.}}
\medskip

Likewise, Arby is willing to sell such a certificate at the same price. Indeed, Arby is willing to buy or sell any number of certificates at this price, as Arby considers it the ``fair'' price.

Arby, not having taken Stat 110, stubbornly refuses to accept the axioms of probability. In particular, suppose that there are two disjoint events $A$ and $B$ with
\[
P_{\text{Arby}}(A\cup B) \ne P_{\text{Arby}}(A) + P_{\text{Arby}}(B).
\]
Show how to make Arby go bankrupt, by giving a list of transactions Arby is willing to make that will guarantee that Arby will lose money (you can assume it will be known whether $A$ occurred and whether $B$ occurred the day after any certificates are bought/sold).
\end{qbox}
\Sol
Identical to Q14. If $P_{\text{Arby}}(A\cup B) < P_{\text{Arby}}(A)+P_{\text{Arby}}(B)$: buy $A$- and $B$-certificates, sell an $(A\cup B)$-certificate, locking in a guaranteed loss of
\begin{align*}
P_{\text{Arby}}(A)+P_{\text{Arby}}(B) - P_{\text{Arby}}(A\cup B) > 0
\end{align*}
per unit traded (disjointness of $A,B$ makes payouts always cancel exactly). If the inequality is reversed, flip all the trades.

%============================================================
\subsection{HW 2 · Problem 2}
\begin{qbox}
A card player is dealt a 13 card hand from a well-shuffled, standard deck of cards. What is the probability that the hand is void in at least one suit (``void in a suit'' means having no cards of that suit)?
\end{qbox}
\Sol
Identical to Q15.
\begin{align*}
P(\text{void in some suit}) = \frac{4\binom{39}{13}-6\binom{26}{13}+4}{\binom{52}{13}} = \frac{4(8122425444)-6(10400600)+4}{635013559600} \approx 0.05107
\end{align*}

\end{document}
